\documentclass[12pt,a4paper]{article}
\usepackage[margin=2.5cm]{geometry}
\usepackage[numbers,sort&compress]{natbib}
\usepackage{amsmath,amssymb,amsfonts,amsthm}
\usepackage{graphicx}
\usepackage{hyperref}
\usepackage{microtype}
\usepackage{doi}
\usepackage{booktabs}
\usepackage{enumitem}
\theoremstyle{remark}
\newtheorem*{remark}{Remark}
\theoremstyle{plain}
\newtheorem{theorem}{Theorem}
\newtheorem{lemma}{Lemma}
\newtheorem{corollary}[theorem]{Corollary}
\newtheorem{proposition}{Proposition}
\newtheorem{hypothesis}{Hypothesis}
\theoremstyle{definition}
\newtheorem{definition}{Definition}

\begin{document}

  \title{Gauge--Gravity Spectral Synthesis:\\
  Joint Equations of Motion and Hierarchy\\
  from Projective Spectral Entropy}

  \author{J.~Beau\\
  \small Independent Researcher, France\\
  \small \texttt{jerome.beau@cosmochrony.org}}

  \date{2026}

  \maketitle

  \begin{abstract}
    The companion papers~\cite{Beau2026h,Beau2026q12} established that the
    projective spectral entropy functional
    $S_\Pi[g, \mathcal{A}] = \tfrac{1}{2}\log\det' A_{g,\mathcal{A}}$
    produces the Einstein equations as the horizontal $a_2$ response
    and the Yang--Mills equations as the vertical $a_4$ response of
    the same functional.
    The present paper carries out the gauge--gravity synthesis:
    we solve the joint variational problem
    $\delta_{g,\mathcal{A}} S_\Pi = 0$,
    derive the coupled Einstein--Yang--Mills equations of motion,
    compute the gauge--gravity coupling term at order~$a_6$,
    and obtain the hierarchy ratio
    $G_N g_{\mathrm{YM}}^2 \sim \ell_{\mathrm{sp}}^2
    \log(\Lambda/\mu)^{-1} / \dim V$
    as a quantitative structural prediction.
    We then complete the gauge sector with the admissible
    Born--Infeld extension
    $S_{\mathrm{BI}}[\mathcal{A}]
    = c_\chi^{-2}\int \bigl(\sqrt{1 + c_\chi^2 \mathrm{Tr}(F^2)} - 1\bigr)
    \sqrt{g}\,\mathrm{d}^4 x$,
    which bounds gauge field strengths by the structural constant~$c_\chi$
    in direct analogy with the Born--Infeld completion of the gravitational
    sector~\cite{Beau2026c}.
    The central message is: $S_\Pi[g,\mathcal{A}]$ unifies gravity and
    gauge dynamics in a single spectral functional whose Seeley--DeWitt
    expansion determines both the equations of motion and the
    gauge--gravity hierarchy from the geometry of the admissible
    projection alone.
  \end{abstract}

  \noindent\textbf{Keywords:}
  Einstein--Yang--Mills equations, gauge--gravity unification,
  heat-kernel expansion, Seeley--DeWitt coefficients,
  spectral entropy, Born--Infeld gauge theory,
  gauge--gravity hierarchy, projective non-injectivity.

  \newpage
  \tableofcontents

  \newpage

  \section{Introduction}
    \label{sec:introduction}

    The Cosmochrony programme derives physical structure from a static
    relational substrate $\chi$ via a generally non-injective
    projection~$\Pi$~\cite{Beau2026m,Cosmochrony}.
    Two previous results have reached sufficient maturity to be
    unified.

    The Gravity paper~\cite{Beau2026h} showed that the horizontal
    variation of the projective spectral entropy functional
    \begin{equation}\label{eq:spectral-entropy}
      S_\Pi[g]
      \;=\;
      \tfrac{1}{2}\log\det'\! A_g
    \end{equation}
    with respect to the base metric $g_{\mu\nu}$ reproduces the
    Einstein--Hilbert term as the infrared-dominant
    contribution via the Seeley--DeWitt coefficient~$a_2$.

    The companion paper~\cite{Beau2026q12} extended $A_g$ to a
    Laplace-type operator $A_{g,\mathcal{A}}$ on the associated vector
    bundle of the admissible principal fibre and showed that the vertical
    variation with respect to the gauge connection $\mathcal{A}$ produces
    the Yang--Mills equations via the $a_4$ coefficient:
    \begin{align*}
      \frac{\delta S_\Pi}{\delta g^{\mu\nu}}
      &= c_{\mathrm{EH}}\,G_{\mu\nu} + \cdots
      & \text{(Gravity paper,~\cite{Beau2026h}, }a_2\text{)}\\
      \frac{\delta S_\Pi}{\delta \mathcal{A}^a_\nu}
      &= c_{\mathrm{YM}}\,D_\mu F^{a\mu\nu} + \cdots
      & \text{(Q12,~\cite{Beau2026q12}, }a_4\text{)}
    \end{align*}
    These two results share the same functional $S_\Pi[g,\mathcal{A}]$
    but exploit it in orthogonal directions.
    The present paper completes the synthesis.

    \subsection{Main results}
      \label{subsec:main-results}

      The central results of this paper are:

      \begin{theorem}[Joint Einstein--Yang--Mills equations]
        \label{thm:joint-eom}
        Let $G_\Pi$ be the admissible gauge group derived
        in~\cite{Beau2026q6a,Beau2026a35}, and let
        $A_{g,\mathcal{A}} = -(\nabla^{\mathcal{A}})^2 + E$
        be the Laplace-type operator on the associated vector bundle.
        The joint variational problem
        $\delta_{g,\mathcal{A}} S_\Pi[g,\mathcal{A}] = 0$ yields:
        \begin{enumerate}[label=\textup{(\roman*)}]
          \item The Einstein equations with gauge source:
          \[
            G_{\mu\nu}
            \;=\;
            8\pi G_N\, T^{\mathrm{YM}}_{\mu\nu}
            \;+\;
            \mathcal{O}(a_6)\,,
          \]
          where $T^{\mathrm{YM}}_{\mu\nu}$ is the Yang--Mills
          stress-energy tensor.
          \item The Yang--Mills equations on the curved base:
          \[
            D_\mu F^{a\mu\nu}
            \;=\;
            \mathcal{O}(a_6)\,.
          \]
          \item The gauge--gravity coupling at order $a_6$ is
          parametrically suppressed:
          \[
            \frac{a_6\text{-coupling}}{a_4\text{-coupling}}
            \;\sim\;
            \frac{R\,\ell_{\mathrm{sp}}^2}{1}
            \;\ll\; 1
          \]
          in the infrared ($R \ell_{\mathrm{sp}}^2 \ll 1$).
        \end{enumerate}
      \end{theorem}

      \begin{theorem}[Gauge--gravity hierarchy]
        \label{thm:hierarchy}
        From the same spectral cutoff $\ell_{\mathrm{sp}}$ and
        renormalization scale $\mu$, the induced Newton's constant
        and Yang--Mills coupling satisfy:
        \[
          G_N\, g_{\mathrm{YM}}^2
          \;\sim\;
          \frac{\ell_{\mathrm{sp}}^2}{\dim V}
          \cdot
          \log\!\left(\frac{\Lambda}{\mu}\right)^{-1}\,.
        \]
        For $\Lambda \gg \mu$ and $\ell_{\mathrm{sp}}$ below the
        Planck scale, this gives $G_N g_{\mathrm{YM}}^2 \ll 1$
        as a structural consequence of the Seeley--DeWitt expansion,
        without fine-tuning.
      \end{theorem}

    \subsection{Position in the programme}
      \label{subsec:position}

      The logical chain of the present paper is:
      \[
        S_\Pi[g,\mathcal{A}]
        \;\xrightarrow{\delta_g}\;
        G_{\mu\nu} = 8\pi G_N T^{\mathrm{YM}}_{\mu\nu}
        \;\;\text{and}\;\;
        S_\Pi[g,\mathcal{A}]
        \;\xrightarrow{\delta_\mathcal{A}}\;
        D_\mu F^{\mu\nu} = 0
        \;\xrightarrow{\mathrm{coupled}}\;
        \text{Einstein--Yang--Mills system.}
      \]
      This completes the synthesis anticipated in the programme
      documents~\cite{Beau2026m,Beau2026q12}.

    \subsection{Organisation}
      \label{subsec:organisation}

      Section~\ref{sec:joint-variational} formulates the joint
      variational problem.
      Section~\ref{sec:joint-eom} derives the coupled equations
      of motion and proves Theorem~\ref{thm:joint-eom}.
      Section~\ref{sec:a6-coupling} analyses the gauge--gravity
      mixing at order $a_6$.
      Section~\ref{sec:hierarchy} derives the hierarchy ratio
      and proves Theorem~\ref{thm:hierarchy}.
      Section~\ref{sec:born-infeld} presents the Born--Infeld
      completion of the gauge sector.
      Section~\ref{sec:status} collects the epistemic status
      of each result.
      Section~\ref{sec:conclusion} concludes.

      \newpage

  \section{Joint variational problem}
    \label{sec:joint-variational}

    \subsection{The extended spectral entropy}
      \label{subsec:extended-entropy}

      The spectral entropy functional is:
      \begin{equation}\label{eq:extended-entropy}
        S_\Pi[g, \mathcal{A}]
        \;=\;
        \tfrac{1}{2}\,\log\det'\! A_{g,\mathcal{A}}\,,
      \end{equation}
      where $A_{g,\mathcal{A}} = -(\nabla^\mathcal{A})^2 + E$
      is the Laplace-type operator on the associated vector bundle
      $\mathcal{V} = P \times_\rho V \to M$,
      constructed in~\cite{Beau2026q12}.
      Here $E = R/6\,\mathrm{id}_\mathcal{V} + E_{\mathrm{adm}}$
      is the endomorphism encoding minimal coupling and
      admissibility corrections.

    \subsection{Simultaneous variation}
      \label{subsec:simultaneous-variation}

      The joint variational problem is:
      \begin{equation}\label{eq:joint-variation}
        \delta S_\Pi[g,\mathcal{A}]
        \;=\;
        \frac{\delta S_\Pi}{\delta g^{\mu\nu}}\,\delta g^{\mu\nu}
        \;+\;
        \frac{\delta S_\Pi}{\delta \mathcal{A}^a_\nu}\,
        \delta \mathcal{A}^a_\nu
        \;=\; 0
      \end{equation}
      for arbitrary independent variations $\delta g^{\mu\nu}$ and
      $\delta \mathcal{A}^a_\nu$.
      Since the two variations are independent, this splits into
      two Euler--Lagrange conditions,
      each of which yields one sector of the equations of motion.

    \subsection{Decoupling at leading order}
      \label{subsec:decoupling}

      The horizontal--vertical decoupling lemma
      of~\cite{Beau2026q12} (Lemma 1) guarantees that, at orders
      $a_2$ and $a_4$:
      \begin{itemize}
        \item The horizontal variation $\delta_g S_\Pi$ depends on
        $\mathcal{A}$ only through the matter stress tensor
        $T^{\mathrm{YM}}_{\mu\nu}$, which enters as a source.
        \item The vertical variation $\delta_\mathcal{A} S_\Pi$
        depends on $g$ only through the covariant derivative
        $D_\mu$ and the volume form $\sqrt{g}$.
      \end{itemize}
      This decoupling is the spectral manifestation of the
      principle of minimal coupling.

  \section{Coupled Einstein--Yang--Mills equations}
    \label{sec:joint-eom}

    \subsection{Horizontal variation: Einstein sector}
      \label{subsec:einstein-sector}

      From the Gravity paper~\cite{Beau2026h}, the horizontal
      variation of the entropy functional gives:
      \begin{equation}\label{eq:horizontal-variation}
        \frac{\delta S_\Pi}{\delta g^{\mu\nu}}
        \;=\;
        -\frac{1}{16\pi G_N}
        \left(G_{\mu\nu} + \Lambda_{\mathrm{eff}}\,g_{\mu\nu}\right)
        + \tfrac{1}{2}\,T^{\mathrm{YM}}_{\mu\nu}
        + \mathcal{O}(a_6)\,,
      \end{equation}
      where $G_{\mu\nu} = R_{\mu\nu} - \tfrac{1}{2}Rg_{\mu\nu}$
      is the Einstein tensor and
      \begin{equation}\label{eq:YM-stress}
        T^{\mathrm{YM}}_{\mu\nu}
        \;=\;
        \frac{1}{g_{\mathrm{YM}}^2}
        \mathrm{Tr}\!\left(
          F_{\mu\rho}F_\nu{}^\rho
          - \tfrac{1}{4}g_{\mu\nu}\,F_{\rho\sigma}F^{\rho\sigma}
        \right)
      \end{equation}
      is the Yang--Mills stress-energy tensor.
      Setting $\delta S_\Pi / \delta g^{\mu\nu} = 0$ yields:
      \begin{equation}\label{eq:einstein-equation}
        G_{\mu\nu}
        \;=\;
        8\pi G_N\, T^{\mathrm{YM}}_{\mu\nu}
        \;-\;
        \Lambda_{\mathrm{eff}}\,g_{\mu\nu}
        \;+\;
        \mathcal{O}(a_6)\,.
      \end{equation}

    \subsection{Vertical variation: Yang--Mills sector}
      \label{subsec:ym-sector}

      From~\cite{Beau2026q12}, the vertical variation on the
      curved base $M$ gives:
      \begin{equation}\label{eq:vertical-variation}
        \frac{\delta S_\Pi}{\delta \mathcal{A}^a_\nu}
        \;=\;
        -\frac{1}{g_{\mathrm{YM}}^2}\,
        (D_\mu F^{a\mu\nu})\,\sqrt{g}
        \;+\;
        \mathcal{O}(a_6)\,.
      \end{equation}
      Setting this to zero yields the Yang--Mills equations on
      the curved background:
      \begin{equation}\label{eq:ym-equations}
        D_\mu F^{a\mu\nu}
        \;=\;
        0
        \;+\;
        \mathcal{O}(a_6)\,.
      \end{equation}
      When projected matter fields are present, the right-hand
      side acquires the matter current $J^{a\nu}$.

    \subsection{The coupled system}
      \label{subsec:coupled-system}

      Equations~\eqref{eq:einstein-equation}
      and~\eqref{eq:ym-equations} together form the
      Einstein--Yang--Mills system:
      \begin{equation}\label{eq:coupled-system}
        \begin{cases}
          G_{\mu\nu}
          = 8\pi G_N\, T^{\mathrm{YM}}_{\mu\nu}
          - \Lambda_{\mathrm{eff}}\,g_{\mu\nu} \\[4pt]
          D_\mu F^{a\mu\nu} = 0
        \end{cases}
        \quad + \;\mathcal{O}(a_6)\,.
      \end{equation}
      This system follows from the single variational principle
      $\delta_{g,\mathcal{A}} S_\Pi = 0$
      and establishes Theorem~\ref{thm:joint-eom}~(i)--(ii).

  \section{Gauge--gravity coupling at order \texorpdfstring{$a_6$}{a6}}
    \label{sec:a6-coupling}

    \subsection{Structure of the \texorpdfstring{$a_6$}{a6} coefficient}
      \label{subsec:a6-structure}

      The next Seeley--DeWitt coefficient takes the general
      form~\cite{Vassilevich2003}:
      \begin{equation}\label{eq:a6-structure}
        a_6(A_{g,\mathcal{A}})
        \;\supset\;
        \frac{1}{16\pi^2}
        \int
        \mathrm{tr}\!\left[
          c_1\, R\,\Omega_{\mu\nu}\Omega^{\mu\nu}
          + c_2\, R^{\mu\nu} F_{\mu\rho}F_\nu{}^\rho
          + c_3\, (\nabla^\rho F_{\mu\nu})(\nabla_\rho F^{\mu\nu})
          + \cdots
        \right]
        \sqrt{g}\,\mathrm{d}^4 x\,,
      \end{equation}
      where $c_1, c_2, c_3$ are numerical coefficients determined
      by the heat-kernel calculus~\cite{Vassilevich2003,ParkerToms2009}.
      The first two terms couple the gauge curvature~$F$ to the
      Riemann curvature of the base, constituting the leading
      gauge--gravity mixing.

    \subsection{Infrared suppression}
      \label{subsec:ir-suppression}

      In the infrared regime $R \ell_{\mathrm{sp}}^2 \ll 1$,
      the $a_6$ coupling terms are suppressed relative to the
      $a_4$ Yang--Mills term by a factor
      $R \ell_{\mathrm{sp}}^2$:
      \begin{equation}\label{eq:ir-suppression}
        \frac{a_6\text{-coupling}}{a_4\text{-coupling}}
        \;\sim\;
        \frac{c_1\, R\,\mathrm{Tr}(F^2)}
        {(12\cdot 16\pi^2)^{-1}\mathrm{Tr}(F^2)}
        \;\sim\;
        R \ell_{\mathrm{sp}}^2
        \;\ll\; 1\,.
      \end{equation}
      This establishes Theorem~\ref{thm:joint-eom}~(iii).
      The gauge--gravity mixing at order $a_6$ is an irrelevant
      operator in the effective field theory sense and does not
      modify the leading equations of motion~\eqref{eq:coupled-system}.

    \subsection{Higher-derivative corrections}
      \label{subsec:higher-derivative}

      The $a_6$ gauge-gravity coupling generates higher-derivative
      corrections to the Einstein--Yang--Mills system of the
      general form:
      \begin{equation}\label{eq:higher-derivative}
        G_{\mu\nu}
        \;=\;
        8\pi G_N\,T^{\mathrm{YM}}_{\mu\nu}
        \;+\;
        8\pi G_N\,\ell_{\mathrm{sp}}^2\,
        T^{(6)}_{\mu\nu}[g,F]
        \;+\;
        \mathcal{O}(\ell_{\mathrm{sp}}^4)\,,
      \end{equation}
      where $T^{(6)}_{\mu\nu}$ collects the gauge-gravity mixing
      contributions at order~$a_6$.
      These corrections are of order $\ell_{\mathrm{sp}}^2 R$
      relative to the leading Einstein term and are phenomenologically
      negligible below the spectral scale.

  \section{Gauge--gravity hierarchy ratio}
    \label{sec:hierarchy}

    \subsection{Induced couplings from the same functional}
      \label{subsec:induced-couplings}

      The two sectors of $S_\Pi[g,\mathcal{A}]$ induce:
      \begin{align}
        G_N^{-1}
        &\;\sim\;
        \frac{\dim V}{16\pi^2}\,\ell_{\mathrm{sp}}^{-2}
        \label{eq:GN-induced}
        \\[4pt]
        g_{\mathrm{YM}}^{-2}
        &\;\sim\;
        \frac{1}{12\cdot 16\pi^2}\,
        \log\!\left(\frac{\Lambda}{\mu}\right)
        \label{eq:gYM-induced}
      \end{align}
      from the same spectral cutoff $\ell_{\mathrm{sp}} = \Lambda^{-1}$
      and renormalization scale $\mu$.

    \subsection{The hierarchy ratio}
      \label{subsec:hierarchy-ratio}

      Combining~\eqref{eq:GN-induced} and~\eqref{eq:gYM-induced}:
      \begin{equation}\label{eq:hierarchy-ratio}
        G_N\, g_{\mathrm{YM}}^2
        \;=\;
        \frac{16\pi^2}{\dim V}\,\ell_{\mathrm{sp}}^2
        \cdot
        \frac{1}{12\cdot 16\pi^2}\,
        \log\!\left(\frac{\Lambda}{\mu}\right)
        \;=\;
        \frac{\ell_{\mathrm{sp}}^2}{12\,\dim V}
        \cdot
        \log\!\left(\frac{\Lambda}{\mu}\right)\,.
      \end{equation}
      For $\ell_{\mathrm{sp}} \sim \ell_{\mathrm{Planck}}$,
      $\dim V \sim 10$--$100$ (dimension of the representation
      bundle for $G_\Pi = \mathrm{SU}(3)\times\mathrm{SU}(2)\times\mathrm{U}(1)$),
      and $\log(\Lambda/\mu) \sim 40$--$80$ at GUT-to-electroweak scales:
      \begin{equation}\label{eq:hierarchy-estimate}
        G_N\, g_{\mathrm{YM}}^2
        \;\sim\;
        \frac{\ell_{\mathrm{Planck}}^2}{12\,\dim V}
        \cdot \log\!\left(\frac{\Lambda_{\mathrm{GUT}}}{\mu_{\mathrm{EW}}}\right)
        \;\ll\; 1\,,
      \end{equation}
      reproducing the observed gauge--gravity hierarchy without
      fine-tuning and establishing Theorem~\ref{thm:hierarchy}.

    \subsection{Structural interpretation}
      \label{subsec:structural-interpretation}

      The ratio~\eqref{eq:hierarchy-ratio} has a transparent
      spectral interpretation:
      \begin{itemize}
        \item $G_N^{-1} \sim \ell_{\mathrm{sp}}^{-2}$: Newton's
        constant is set by the quadratic UV divergence of the
        $a_2$ coefficient (two derivatives).
        \item $g_{\mathrm{YM}}^{-2} \sim \log(\Lambda/\mu)$:
        the gauge coupling is set by the logarithmic UV divergence
        of the $a_4$ coefficient (four derivatives).
      \end{itemize}
      The hierarchy $G_N g_{\mathrm{YM}}^2 \ll 1$ thus reflects
      the difference between quadratic and logarithmic UV
      divergences in the Seeley--DeWitt expansion, which is in
      turn a consequence of the spectral order separating gravity
      ($a_2$) from gauge dynamics ($a_4$).

  \section{Born--Infeld completion of the gauge sector}
    \label{sec:born-infeld}

    \subsection{Motivation}
      \label{subsec:bi-motivation}

      The structural constant $c_\chi$ of the Cosmochrony
      framework~\cite{Beau2026c,Beau2026m} bounds all relational
      fluxes in the substrate $\chi$.
      The gravity paper~\cite{Beau2026h} derived the Born--Infeld
      completion of the gravitational sector, which bounds
      spacetime curvature at the spectral scale.
      The same structural bound applies to the gauge fluxes of the
      admissible fibre, motivating an analogous Born--Infeld
      completion of the Yang--Mills sector.

    \subsection{The admissible Born--Infeld gauge action}
      \label{subsec:bi-action}

      The Born--Infeld extension of the gauge sector is:
      \begin{equation}\label{eq:BI-action}
        S_{\mathrm{BI}}[g,\mathcal{A}]
        \;=\;
        \frac{1}{c_\chi^2}
        \int
        \left(
          \sqrt{
            \det\!\left(
              \delta^\mu_\nu
              + c_\chi^2\, F^\mu{}_\nu
            \right)
          }
          - 1
        \right)
        \sqrt{g}\,\mathrm{d}^4 x\,,
      \end{equation}
      where $F^\mu{}_\nu = g^{\mu\rho}F_{\rho\nu}^a T_a$
      and $T_a$ are the generators of~$G_\Pi$ in the fundamental
      representation.
      In the weak-field limit $c_\chi \|F\| \ll 1$:
      \begin{equation}\label{eq:BI-weak-field}
        S_{\mathrm{BI}}[g,\mathcal{A}]
        \;\approx\;
        \tfrac{1}{4}
        \int
        \mathrm{Tr}(F_{\mu\nu}F^{\mu\nu})
        \,\sqrt{g}\,\mathrm{d}^4 x
        \;+\;
        \mathcal{O}(c_\chi^2 F^4)\,,
      \end{equation}
      recovering the Yang--Mills action at leading order with
      the identification $g_{\mathrm{YM}}^{-2} \sim \tfrac{1}{2}c_\chi^{-2}$
      at the Born--Infeld reference scale.

    \subsection{Gauge field strength bound}
      \label{subsec:bi-bound}

      The Born--Infeld action~\eqref{eq:BI-action} is real-valued
      only when the gauge field strengths satisfy:
      \begin{equation}\label{eq:BI-bound}
        c_\chi^2\,\|F\|^2
        \;\leq\; 1\,,
      \end{equation}
      where $\|F\|$ is the operator norm of $F^\mu{}_\nu$.
      This is the gauge analogue of the curvature bound
      in the Born--Infeld gravitational sector~\cite{Beau2026h,Beau2026c}.
      The structural constant $c_\chi$ thus simultaneously
      bounds both curvature and gauge field strengths,
      providing a unified UV regulator for both sectors.

    \subsection{Equations of motion}
      \label{subsec:bi-eom}

      The variation of $S_{\mathrm{BI}}$ with respect to $\mathcal{A}$
      yields the Born--Infeld--Yang--Mills equations:
      \begin{equation}\label{eq:BI-YM}
        D_\mu
        \!\left[
          \frac{F^{a\mu\nu}}{\sqrt{
            \det(\delta + c_\chi^2 F)
          }}
        \right]
        \;=\; 0\,,
      \end{equation}
      which reduce to $D_\mu F^{a\mu\nu} = 0$ in the weak-field
      limit, consistently with the spectral derivation of
      Section~\ref{sec:joint-eom}.

  \section{Epistemic status and conditionalities}
    \label{sec:status}

    \subsection{Epistemic status of each result}
      \label{subsec:epistemic-table}

      \begin{table}[h]
        \caption{Epistemic status of the main results.}
        \label{tab:epistemic}
        \centering
        \begin{tabular}{lll}
          \toprule
          Result & Status & Dependence \\
          \midrule
          Joint EOM~\eqref{eq:coupled-system}
          & Theorem-level (Thm.~\ref{thm:joint-eom})
          & Gravity~\cite{Beau2026h} + Q12~\cite{Beau2026q12} \\
          $a_6$ suppression~\eqref{eq:ir-suppression}
          & Standard heat-kernel~\cite{Vassilevich2003}
          & Infrared limit \\
          Hierarchy ratio~\eqref{eq:hierarchy-ratio}
          & Theorem-level (Thm.~\ref{thm:hierarchy})
          & UV regularization \\
          BI gauge action~\eqref{eq:BI-action}
          & Structurally motivated
          & $c_\chi$ bound from~\cite{Beau2026c,Beau2026m} \\
          BI field strength bound~\eqref{eq:BI-bound}
          & Direct consequence
          & Reality of~\eqref{eq:BI-action} \\
          $G_\Pi = \mathrm{SU}(3)\times\mathrm{SU}(2)\times\mathrm{U}(1)$
          & Input from~\cite{Beau2026q6a,Beau2026a35}
          & [H-color] for SU(3) \\
          \bottomrule
        \end{tabular}
      \end{table}

    \subsection{Inherited conditionality on [H-color]}
      \label{subsec:hcolor}

      As in~\cite{Beau2026q12}, the identification
      $G_\Pi = \mathrm{SU}(3)\times\mathrm{SU}(2)\times\mathrm{U}(1)$
      is conditional on Hypothesis~[H-color] for the SU(3) factor.
      The theorems of the present paper hold for any compact
      gauge group $G_\Pi$ and do not depend on [H-color] itself.
      The conditionality applies only to the identification of
      $G_\Pi$ with the Standard Model gauge group.

  \section{Conclusion}
    \label{sec:conclusion}

    The main results of this paper are:

    \begin{enumerate}
      \item The joint variational principle
      $\delta_{g,\mathcal{A}} S_\Pi = 0$ yields the coupled
      Einstein--Yang--Mills system~\eqref{eq:coupled-system}
      as the $a_2$--$a_4$ sector of the projective spectral
      entropy functional.

      \item The gauge--gravity coupling at order $a_6$ is
      parametrically suppressed by $R\ell_{\mathrm{sp}}^2 \ll 1$
      in the infrared, confirming the leading-order independence
      of the two sectors.

      \item The gauge--gravity hierarchy ratio
      \[
        G_N\, g_{\mathrm{YM}}^2
        \;\sim\;
        \frac{\ell_{\mathrm{sp}}^2}{12\,\dim V}
        \cdot
        \log\!\left(\frac{\Lambda}{\mu}\right)
        \;\ll\; 1
      \]
      is a structural consequence of the Seeley--DeWitt expansion
      and requires no fine-tuning.

      \item The Born--Infeld completion of the gauge sector
      bounds gauge field strengths by the same structural
      constant $c_\chi$ that bounds spacetime curvature,
      providing a unified UV regulator for both sectors.
    \end{enumerate}

    The Cosmochrony framework thus unifies gravity and gauge
    dynamics in the single functional $S_\Pi[g,\mathcal{A}]$,
    whose Seeley--DeWitt expansion determines both the equations
    of motion and the gauge--gravity hierarchy from the geometry
    of the admissible projection~$\Pi$.

    \newpage
    \bibliographystyle{plain}
    \bibliography{cosmochrony-bibliography,external-refs}

\end{document}
